\section{What the five global estimates must prove, and two sufficient simplifications}
\label{sec:global-quantitative-v042}

\status{This chapter proves the numerical and finite-counting passage to
the original retained-row defect, on the exact stated row inputs. It
also proves that a uniform polynomial spectral error suffices for the
selected final comparison. The actual arithmetic row counts, local
weights, period estimates and spectral estimates remain input-producing
theorems. None of the five whole global obligations is declared closed.}

The source is the pinned revision
\texttt{7cbf7a0c32fca1badf05931879a37c3f8a2b03b4}.
Every statement below concerns the original family and its original
maps. In particular, the letters denoting source spaces, target spaces,
Haar measures, metrics, coefficient embeddings and auxiliary primes
do not stand for independently chosen models. The source correspondence
and full declarations are recorded in the companion
\texttt{contracts.json} and \texttt{source\_manifest.json}.

\subsection{The five estimates as mathematical contracts}

Fix the family, its positive small exponent $\epsilon$, and its height
exponent $C_H$. At scale $x$ let $D$ denote any actual member, and put
\[
 L=x^\epsilon,\qquad b=e^L,\qquad a=xb,\qquad
 B_0=fmbL,\qquad \Gamma=e^{C_\Gamma m^5fL}.
\]
The thresholds in this subsection are chosen before the member, word,
lift, or finite place. A statement described as closed over a class
additionally preserves the source's earlier order of choosing the
small-exponent cutoff and the parameter-bound functions.

\paragraph{Denominators.}
For a lift $j$, let $d_{Z,0}(j)$ be the actual small-level rank and set
\[
 \mathcal E_j=\sum_{\substack{w\in\mathrm{supp}\\\mu(w)\ne\mu_0}}
       \ell_w e_w\eta_j(w),\qquad
 V(D)=\sum_{p\ \mathrm{prime}}\log p
                 \sum_{1\le e<t}\operatorname{shaLen}(D,p,e).
\]
Here $\ell_w$ is exactly the source's logarithmic place weight, $e_w$
its ramification index, and the finite cohomology lengths are read from
the same arithmetic datum $D.\mathrm{nd}$. Under the original
finite-length and summability guards, the required estimate is
\begin{equation}\label{eq:gqt-denominator}
 \mathcal E_j\le d_{Z,0}(j)
   \left[C\{fmbL^2+(\log x)^2N_c\}+V(D)\right].
\end{equation}
One $C$ works for all late members and all lifts. This is
\texttt{LinkN}, including its guards; a bound using a different
cohomology datum is not the same assertion.

\paragraph{Retained rows.}
Write $D_{\rm all}$ and $D_{\rm rows}$ for the actual full and retained
target dimensions. The stored defect satisfies
\[
 D_{\rm rows}=(1-\delta)D_{\rm all}.
\]
For $\epsilon_z=\epsilon/J>0$, the required assertion is
\begin{equation}\label{eq:gqt-defect}
                  0\le\delta\le e^{-x^{\epsilon_z/7}}.
\end{equation}
The source's theorem quantifies over every late member. The proofs
below isolate exactly the counting and probability estimates needed
for this conclusion.

\paragraph{Periods and heights.}
Let $R$ be the common retained target lattice, $N_u$ the actual shifted
source lattice, $B_u:N_u\to R^\vee$ the rational period map, and
$h_u$ its supplied squared-norm height majorant. Besides integrality,
rationality, fullness of $R$, the support cutoff and the exact Bessel
identity, \texttt{LinkDIV} requires
\begin{align}
 s_{\rm Frame}+\mu_{\max}(R)&\le x^{C_H}b,
                                                        \label{eq:gqt-height1}\\
 h(B_u)&\le h_u\quad(B_u\ne0),                           \label{eq:gqt-height2}\\
 \mu(R^\vee)+h_u-\mathrm{Fo}
   &\le\mathrm{MAIN}_u+C n_TaL_a
       +Cm(B_0+N_c)x^{C_{\rm abs}\epsilon},               \label{eq:gqt-height3}\\
 \left|\sum_w\mathrm{sz}(w,E_{u,j,w})
                  -A_T\log(g_0g_1)\right|&\le CB_0,     \label{eq:gqt-height4}\\
 \left|(h_u+A_T\log T_u)-(h_v+A_T\log T_v)\right|
       &\le C(n_TaL_a+mB_0).                             \label{eq:gqt-height5}
\end{align}
Here $n_T=\#TT$, $A_T=\sum_{t\in TT}a_t$, and every term on
the right is the original family quantity. In particular
\[
 \begin{aligned}
 \mathrm{MAIN}_u&=
 \sum_{t\in TT}\bigl[m(\phi(a_t)-a_t\log\Omega)-a_t\log T_u\bigr],\\
 \phi(y)&=y\log y-y,\\
 \log\Omega&=\gamma+\log(\pi/v)-L'/L(1).
 \end{aligned}
\]
The Bessel identity uses one target isometry before all words, then a
nonzero scalar and a source isometry for each word. The exported
period-height implication has the precise order
\[
 \mathrm{TuSize}\Longrightarrow\ \mathrm{DGamSize}\
 \Longrightarrow\bigl(\mathrm{LinkEta}\Longrightarrow
                     \mathrm{DIVinN}\Longrightarrow\mathrm{LinkDIV}\bigr).
\]
It does not assert that those upstream inputs have already been
constructed.

\paragraph{Local gains.}
At places with residue characteristic different from $T_u$, the
original local lattices must satisfy the stated valuation inclusion,
\[
 \operatorname{len}(M/N)_{u,w}\le d_u^{\rm full}
                                      (\tau_w+\mathrm{io}_w),
\]
and $\tau_w=\mathrm{io}_w=\operatorname{len}(M/N)_{u,w}=0$
when there is no served orbit. A served orbit $t_0$ at $p$ must obey
$S_v(p,t_0)\le2mN_v(p,t_0)a/p$. At the specified signal or
inert-class primes $p\le X''$, one also needs
\begin{gather}
       0\le\mathrm{gain}_{j,w}\le4n_Ta/p,                 \label{eq:gqt-gain1}\\
 \sum_j\sum_{t\in TT_w}
  \{\mathrm{vonMangoldtClassSum}_{j,t,p}-\mathrm{loG}(C,b_o,p)\}
       -m\,\mathrm{flG}(C,p)
   \le\sum_j\mathrm{gain}_{j,w}-\tau_w.                 \label{eq:gqt-gain2}
\end{gather}
The same assertion includes served-orbit containment and
\texttt{MeanKappaBound}. The full formulas for the two error
envelopes are retained verbatim in the source contract record;
in particular $\mathrm{flG}(C,p)=CfN(m)^2\log_p a+
Cf(\log_p a)^2$. There must be a single function $\rho(x)\to0$
with $b_o=b^{\rho(x)}$ before the family index is varied.
The source does not posit separately chosen losses for each prime.

\paragraph{Spectrum.}
For the full period operators set $A_u=c_u^{-1/2}P_u$ using
\[
 c_u=\frac{\operatorname{vol}(N_u^{\rm an})
                 \operatorname{vol}(H_u^{\rm an})}{\dim\Pi_u}.
\]
The original proviso is
\begin{equation}\label{eq:gqt-proviso}
 C_0^{-1}\le d_u\Gamma^2/D_*\le C_0\qquad\text{for every }u.
\end{equation}
Conditional on this proviso, the desired normalized inequalities are
\begin{align}
 \|A_u^*A_v\|_{\mathrm{HS}}^2
     &\le(1+\eta_0)d_ud_v/D_*\quad(u\ne v),             \label{eq:gqt-off}\\
 \|A_u^*A_u-I\|_{\mathrm{HS}}^2&\le e_0d_u,            \label{eq:gqt-diag}\\
 D_{\rm all}&\le(1+\eta_1)D_*.                         \label{eq:gqt-dimension}
\end{align}
The full original \texttt{LinkS} asks for $e_0\ge0$ and
$\eta_0,\eta_1,\sqrt{e_0}\le e^{-x^{\epsilon/8}}$.
The errors may depend on the member, but their displayed upper bound
and late threshold must be uniform. GQT6--GQT7 below prove a weaker
sufficient error rate for the selected final comparison. They do not
produce the period operators or prove these estimates.

\subsection{Only three counts are needed for the retained target}

% BEGIN GQT1
\paragraph{GQT1: a finite union estimate with relative counting errors.}
\label{claim:GQT1}
Let $T$ be a finite set of $n$ rows, with two subsets $B_0,B_1$ of
bad rows. Assume that the good rows are exactly
$G=T\setminus(B_0\cup B_1)$, and write $g=\#G$.
Suppose $M>0$, $0\le\alpha<1$, and $p_0,p_1\ge0$ satisfy
\[
 |n-M|\le\alpha M,\qquad
 |\#B_j-p_jM|\le\alpha p_jM\quad(j=0,1).
\]
Then $n>0$, and for any $W_j\ge p_j$,
\begin{equation}\label{eq:gqt-three-counts}
 0\le1-\frac gn
       \le\frac{1+\alpha}{1-\alpha}(p_0+p_1)
       \le\frac{1+\alpha}{1-\alpha}(W_0+W_1).
\end{equation}
Indeed $n\ge(1-\alpha)M>0$. The complement of $G$ is the union
of the two bad subsets, so $n-g\le\#B_0+\#B_1
\le(1+\alpha)(p_0+p_1)M$. Divide by the lower bound for $n$.
Finally $g\le n$ proves nonnegativity. No independence of the two
bad events is required.

In fact only the lower bound for $n$ and upper bounds for the two
bad counts were used. The opposite halves of the three absolute-value
estimates, and counts for every other subset, may be omitted.
% END GQT1

% BEGIN GQT2
\paragraph{GQT2: these are the original total and bad-cylinder counts.}
\label{claim:GQT2}
Supply one original \texttt{AllRows} object $\mathcal A$ for $D$,
together with \texttt{LabelCountGoodInjective}. Thus its all-row
labels $\lambda_j(\tau)$ are indexed by
$\tau\in T=\{1,\ldots,D_{\rm all}\}$; its injection of the
$D_{\rm rows}$ retained labels has image exactly the labels good
at both auxiliary places. Its dimension function counts these
labels in every cylinder. For $j=0,1$, put
\[
 B_j=\{\tau\in T:\lambda_j(\tau)\text{ is not good}\},
 \qquad S_j=\{\lambda_j(\tau):\tau\in B_j\}.
\]
These are finite sets even though the ambient polynomial-label space
is infinite. A row lies in the cylinder determined by $S_j$ exactly
when it belongs to $B_j$: every label in $S_j$ is bad, and each bad
row contributes its label. Consequently
\[
 \dim\mathcal A(\mathrm{cyl}(j,S_j))=\#B_j,
 \quad \dim\mathcal A(\mathrm{cyl}(0,\mathrm{all}))=D_{\rm all},
 \quad\#(T\setminus(B_0\cup B_1))=D_{\rm rows}.
\]
The last identity uses the supplied injection and its exact
good-image assertion, not an assumption about multiplicities of
distinct labels.

Let
\[
 p_j=\sum_{t\in S_j}\prod_v
   \frac{\operatorname{Re}(e_{j,v}e_{j,t(v)})(1)}
        {\operatorname{Re}e_{j,v}(1)}
       =\operatorname{locShare}(\mathcal A,j,S_j),
 \qquad M=D_*^{\mathcal A}(\mathcal A.B.e).
\]
The notation on the right fixes the source convention even in any
zero-denominator branch. Supply the original nonnegative local
weight assertion, so $p_j\ge0$. The source's finite cylinder
identity states
$\operatorname{Re}e_{\mathrm{cyl}(j,S_j)}(1)
=\operatorname{Re}e(1)p_j$. Its definition
$D_*^{\mathcal A}(e)=d_{\rm formal}\operatorname{vol}\,
\operatorname{Re}e(1)$ therefore gives
\[
       D_*^{\mathcal A}(e_{\mathrm{cyl}(j,S_j)})=Mp_j.
\]
No change of Haar measure or formal degree has entered this equality.

The original \texttt{CountAllRows} supplies the three relative
estimates in GQT1 and $M>0$. The weight predicate supplies
$p_j\le W_g(f,m,N,g_j,q,d)$, with $N=D.R.r$ and $g_j=D.host.g(j)$.
The stored identity $D_{\rm rows}=(1-\delta)D_{\rm all}$ now
identifies the defect in GQT1 with the literal $D.\mathrm{what}_2.
\mathrm{delp}$. Thus
\begin{equation}\label{eq:gqt-source-share}
 0\le\delta\le\frac{1+\alpha}{1-\alpha}
       \bigl(W_g(f,m,N,g_0,q,d)+W_g(f,m,N,g_1,q,d)\bigr).
\end{equation}
This proof needs the counting estimate only for the full set and the
two finite bad cylinders selected above. It needs the weight bound
only for these two sets. These weaker hypotheses therefore give a
proved alternative input contract for the original defect consumer.
They do not assert that these three arithmetic estimates are known.
% END GQT2

\subsection{The uniform numerical estimate is elementary}

The explicit source envelope used in GQT2 is
\begin{align}
 W_g(f,m,N,g,q,d)={}&\frac{3f}{g}
       +\frac{23f\log g}{\sqrt g}
       +2fm\left\{\frac1{4\min(q,d)}
                         +\frac{1+\log N}{g-1}\right\}\notag\\
       &+\frac N2(7/8)^f
         +\exp\left(-\frac f{5\log N}\right).           \label{eq:gqt-W}
\end{align}

% BEGIN GQT3
\paragraph{GQT3: one threshold for every admissible numerical tuple.}
\label{claim:GQT3}
Fix $0<\vartheta\le1$ and $C_s\in\mathbb R$. For all sufficiently
large $x$, simultaneously for every integer tuple satisfying
\begin{equation}\label{eq:gqt-sizes}
 g\ge e^{8x^\vartheta},\quad x^2\le f\le x^{C_s},\quad
 m\le x^{C_s},\quad100\le N\le x^{C_s},\quad
 \min(q,d)\ge e^{x^{\vartheta/6}},
\end{equation}
one has
\begin{equation}\label{eq:gqt-Wbound}
         W_g(f,m,N,g,q,d)
       \le x^{3C_s+3}e^{-x^{\vartheta/6}}.
\end{equation}
The threshold depends only on $\vartheta,C_s$, not the tuple.

If $C_s<2$, there are no admissible tuples for $x>1$. Assume
$C_s\ge2$, and abbreviate $t=x^\vartheta$, $v=x^{\vartheta/6}$.
For late $x$, $g\ge2$ and $\log g\ge8t\ge2$.
The function $z e^{-z/2}$ is decreasing for $z\ge2$, whence
\[
 \frac{\log g}{\sqrt g}\le8t e^{-4t},\qquad
 \frac1{g-1}\le2e^{-8t}.
\]
Use these inequalities, $\log N\le C_s\log x$,
and $(7/8)^f=e^{-f\log(8/7)}$ in \eqref{eq:gqt-W}. This gives
\begin{align*}
 W_g\le{}&3x^{C_s}e^{-8t}
       +184x^{C_s+\vartheta}e^{-4t}
       +\tfrac12x^{2C_s}e^{-v}
       +4x^{2C_s}(1+C_s\log x)e^{-8t}\\
       &+\tfrac12x^{C_s}e^{-\log(8/7)x^2}
       +\exp\left(-\frac{x^2}{5C_s\log x}\right).
\end{align*}
After multiplication by $e^v$, the first, second, fourth, fifth
and sixth terms tend to zero. For the first three of these, use
$v=o(t)$ and $\log x=o(t)$; for the fifth use $v=o(x^2)$;
for the sixth use $v\log x=o(x^2)$, valid since
$\vartheta/6\le1/6$. These statements follow, for example, by
writing $x=e^y$ and comparing a fixed power of $y$ with
$e^{cy}$ for $c>0$, using its power series. They do not conceal
constants depending on $f,m,N,g,q,d$.

Choose one threshold at which their sum is at most one. Then
$W_g\le(1+x^{2C_s}/2)e^{-v}$, which is at most
$x^{3C_s+3}e^{-v}$ for $x\ge2$. This proves
\eqref{eq:gqt-Wbound} with the source's stated exponent.
% END GQT3

% BEGIN GQT4
\paragraph{GQT4: absorb the counting error and obtain the exact defect.}
\label{claim:GQT4}
Fix also $c>0$ and set
\[
                 \alpha=e^{-cNf\log N}.
\]
Under \eqref{eq:gqt-sizes},
$\alpha\le e^{-100c\log(100)x^2}$, so $0<\alpha\le1/2$
uniformly for late $x$. Therefore
$(1+\alpha)/(1-\alpha)\le3$. Applying GQT3 at both $g_0,g_1$
in \eqref{eq:gqt-source-share} gives
\[
          0\le\delta\le6x^{3C_s+3}e^{-x^{\vartheta/6}}.
\]
Finally
\[
 x^{\vartheta/6}-x^{\vartheta/7}
                   -(3C_s+3)\log x-\log6\longrightarrow+\infty,
\]
because $x^{\vartheta/7}=o(x^{\vartheta/6})$ and
$\log x=o(x^{\vartheta/6})$. Enlarge the common threshold once.
Then
\begin{equation}\label{eq:gqt-final-defect}
                         0\le\delta\le e^{-x^{\vartheta/7}}.
\end{equation}
Every threshold in this argument depends only on $\vartheta,C_s,c$.
It precedes all row data, labels and numerical tuples.
% END GQT4

% BEGIN GQT5
\paragraph{GQT5: the exact conditional family endpoint, and its smaller input.}
\label{claim:GQT5}
Fix $c>0$, $C_s\in\mathbb R$, and the source family $\mathcal F$
with $0<\epsilon\le1$,
$J\ge1$, and $\epsilon_z=\epsilon/J$. Suppose that for every late
member there is one original $\mathcal A:\texttt{AllRows}(D)$ with
\texttt{LabelCountGoodInjective}, \texttt{CountAllRows} at
$\alpha=e^{-c\,D.R.r\,D.R.f\log(D.R.r)}$, and
\texttt{WeightAllRows}; also suppose \texttt{SizesRec} at
$(x,\epsilon_z,C_s)$. Then the literal \texttt{LinkIRR} holds.

Indeed $0<\epsilon_z\le1$. Take the maximum of the supplied
family threshold and the thresholds of GQT3--GQT4 for
$\vartheta=\epsilon_z$. At any later $x$, choose the supplied
row object for each member and apply GQT2--GQT4. The resulting
\eqref{eq:gqt-final-defect} is exactly the defect and exponent
in \eqref{eq:gqt-defect}. This proves the mathematical content
of \texttt{ConcN.linkIRR\_of\_allRows}, including positivity,
the original row dimension, and the uniform quantifiers.

The same proof works if the full cylinder-count and weight hypotheses
are replaced by the three inequalities and two finite-set weight
bounds specified in GQT2. As an even weaker sufficient version, one
may take any $0\le\alpha\le1/2$ with those three inequalities:
the particular exponential formula for $\alpha$ is not required
after that uniform bound is established. The source's strict rate
for $\delta$ still follows from the same $W_g$ estimates.

This closes the numerical implication, not the global existence
and counting problem. In particular the original all-row object,
its label/good-image identification, and its local weights are
not inferred from a list of dimensions. The separate source route
\texttt{AllJoinA}\,$\to$\,\texttt{WeightAllRows} retains its
finite-type and character calculations.
% END GQT5

\subsection{The selected comparison needs a polynomial spectral error}

% BEGIN GQT6
\paragraph{GQT6: a weaker spectral rate gives the full frame conclusion.}
\label{claim:GQT6}
Fix $C_H>0$, $C_0\ge1$, and put $\sigma=x^{-C_H}$. Retain the actual normalized
operators and the structural hypotheses of GFA6--GFA8: the original
proviso \eqref{eq:gqt-proviso}, code cardinality
$\#W\ge C_0\Gamma^2/\sigma$, ramified source inclusions,
\texttt{LinkIRR}, and the full \texttt{LinkDIV}. Retain also the
full positive-rank source lattices and their original rank, height,
and retained-subspace identifications. In place of
\texttt{LinkS}, suppose that
\eqref{eq:gqt-off}--\eqref{eq:gqt-dimension} hold at every late
member with $e_0\ge0$ and
\begin{equation}\label{eq:gqt-poly}
       \eta_0,\eta_1,\sqrt{e_0}\le\sigma/16.
\end{equation}
Negative $\eta_0,\eta_1$ may be replaced by zero. Suppose also
the actual discarded source fraction is at most $2f/b$, with
$f\le x^5$ and $b=e^{x^\epsilon}$ as in GFI15.
Then the same full \texttt{FrameF}$(\mathcal F,C_H)$ holds.

Here is the complete error-budget check. Choose one threshold with
\[
 \sigma\le1/20,\qquad 2x^5e^{-x^\epsilon}
                                      \le\sigma^2/256.
\]
The second follows by taking logarithms, since
$x^\epsilon-(5+2C_H)\log x-\log512\to+\infty$.
Thus we may use
\[
      d=\sigma^2/256,\qquad
      \overline\eta=83\sigma/16,
      \overline e=\sigma^2/128.                         \label{eq:gqt-frame-errors}
\]
The proviso gives $d_u>0$. The diagonal error has $e_0<1$,
so the same argument as GFA6 gives $D_{\rm all}>0$; otherwise
the period operators would be zero and the diagonal squared
defect would equal $d_u$. Likewise the nonnegative normalizing
constant $c_u$ cannot be zero. The original row defect is less
than one, so fullness and the retained-dimension identity give
$\operatorname{rank}R>0$. These checks use no exponential
spectral rate. The coefficient and Bessel-map identifications
in GFA7 therefore apply unchanged.

Now $\eta_0,\eta_1,\sqrt{e_0}\le\sigma/2$ and
$d\le1/4$. Apply the proved Gram-to-slope calculation GF8,
or equivalently its explicit specialization GFA8. For every
$H>0$ with $s-\mu_{\min}(R^\vee)\le H$, it gives the actual
frame inequality
\[
 s-\log\#W-
 \frac{4u\log(1+H/u)+(\eta+4\sqrt e)H}{1-\delta}
                         \le\mu(R^\vee),
\]
with $\eta=5\sigma+3\sqrt d$ and $e=2e_0$.
The term $u\log(1+H/u)$ is interpreted as zero at $u=0$,
as in the inherited statement. The chosen errors obey
\[
 0\le\eta\le\overline\eta,\quad
 0\le e\le\overline e\le1/4,
 \qquad
 (\overline\eta+4\sqrt{\overline e})x^{C_H}
       =\frac{83+4\sqrt2}{16}<6.
\]
These are the literal all-height inequality and uniform error
conditions required in \texttt{FrameF}. The errors
$\overline\eta,\overline e$ are chosen before the family member;
the original $\delta$ remains distinct from the discarded source
fraction $d$. This proves the claim.
% END GQT6

% BEGIN GQT7
\paragraph{GQT7: the weaker spectral estimate is sufficient for the final consumer.}
\label{claim:GQT7}
Keep every hypothesis of GFI15 except its original exponential
\texttt{LinkS}. Replace that hypothesis by
\eqref{eq:gqt-off}--\eqref{eq:gqt-dimension} and
\eqref{eq:gqt-poly}, with one threshold before all members at
the fixed $C_H$. Then the same numerical conclusion appearing after
the antecedents of \texttt{LinkGLOBAL7} follows with
$e_D=C_{\rm abs}$ and $e_L=0$:
\[
 \exists K,c_{HH},e_K,C_L,x_0\quad
 \forall x\ge x_0,\ x\in\mathcal F.\mathrm{Win}\ \Longrightarrow\
 \mathrm{ChainTied}(\mathcal F,C_{\rm abs},0,K,c_{HH},e_K,C_L,x).
\]

To check every use, GFI15 first constructs the actual ramified
dimension formula and source-lattice fullness from its supplied
local fields and proviso. Its counting, source-slope,
quotient-slope, local-loss, node and character arguments do not
use \texttt{LinkS}. The remaining uses are GFA6's positive
target rank, GFA7's exact normalization and realization, and
GFA8--GFA9's frame inequality and error budget. GQT6 supplies
all of these under the replacement. Thus GFI6--GFI13 receive
the same nine family predicates, the same lattices, the same
height exponent, and the same bound
$(\overline\eta+4\sqrt{\overline e})x^{C_H}\le6$.
Their numerical calculation is unchanged, which proves the
displayed conclusion without assuming the stronger spectral predicate.
The subsequent already
proved window argument retains exponent $1/2+9\epsilon$.

For a closed family-class assertion, the replacement itself must
be supplied with the source's earlier parameter-choice order and
uniformity. A separately chosen bound after fixing each member
would not suffice. For the human proof at fixed $C_H$, any
uniform estimate
\[
 \eta_0,\eta_1,\sqrt{e_0}=O(x^{-C_H-\rho})\qquad(\rho>0)
\]
implies \eqref{eq:gqt-poly} by increasing the threshold.
The extra exponent need not be specified in advance if the
uniform little-oh estimate $o(x^{-C_H})$ is proved.

This is a replacement of the spectral \emph{input to the selected
consumer}. It does not prove the original stronger
\texttt{LinkS}, does not identify every other use of that
predicate in the full Lean development, and does not eliminate
the actual spectral-geometry work. The unaltered global
obligation remains open until either the original estimates or
this sufficient alternative have been constructed on the common
family with all the stated normalizations.
% END GQT7

\subsection{Reconstruction tasks exposed by the exact contracts}

For the denominator estimate, the next substantive target is the
actual local scalar table and its common choice of bounds. The
source passes from \texttt{NormChoiceFacts}, the uniform
\texttt{Extra} bounds and the two/three-place bounds to
\texttt{ConclF}, then to \texttt{AdmF}, then to
\eqref{eq:gqt-denominator}. The equality between the exported
$\eta$-sum and the source sum is definitional. Proving that equality
again would not produce the missing estimates. The Galois cut
transport must preserve the original row rank and multiplicity
bound. The guards concerning the actual cohomology lengths remain.

For the retained target, GQT1--GQT5 remove the finite counting
and asymptotic tail from the reconstruction frontier. The
remaining work is to construct the actual labels and common
row object and prove the three required count estimates and two
bad-weight bounds. The full \texttt{CountAllRows} and
\texttt{WeightAllRows} producers remain sufficient; the weaker
contracts offer an alternative if they reduce the arithmetic work.

For periods and heights, a useful division is: actual rational
common target and exact Bessel map; all-embedding chain norms and
eigenvalue lower bounds; local divisor/arithmetic comparison;
the two finite-place and cross-word remainder estimates; and
the uniform height bound \eqref{eq:gqt-height1}. Each must refer
to the same $R,N_u,B_u$ and the same metrics. The existing
period-lattice and metric-transport arguments should be reused.
The source's \texttt{RemainingDivisionBounds} is only part of
the full contract, not a synonym for \texttt{LinkDIV}.

For local gain, the source reduces the global installation to
the same \texttt{GroupCompatibility\allowbreak{}Family} and admissible-top
condition, then actual local extremal-projector, lattice, flatness
and descent estimates. The \texttt{Loc} parameter supplies the
uniform \texttt{SecondTheoremReduced}; the local-model and
flatness constructors use further fixed formal theorems. Each
relevant result still needs a human proof or a checked published
application. The uniform subexponential
$b^{\rho(x)}$ loss and the original orbit representatives are
part of its output.

For the spectrum, retain separately the exact normalized period
expansion; the degenerate, nongeneric, loop and generic terms;
the Fourier/derivative and lattice-summation argument; and the
independent target-dimension estimate. The original source's
\texttt{TheoremSSelected}, \texttt{LoopRowS},
\texttt{RowNGP}, \texttt{RowMainEqP}, \texttt{TypeP},
\texttt{WuDecayOfGap}, \texttt{RowsFinDim}, and all-point count
are genuine input theorems. The earlier AG chapter already
proves the finite recurrence and a sufficient packing estimate.
GQT6--GQT7 lower the required final accuracy; they do not replace
these constructions by formal theorem names.

The new mathematics is therefore a complete conditional
retained-defect theorem, a smaller sufficient row-count contract,
and a proved weaker spectral input for the selected final
comparison. The common-family and arithmetic existence
obligations retain their previous status.
